% ===== §4 The Second Cause: Unattributable Prediction Error =====
\section{The Second Cause: The Unattributability of Prediction Error}
\label{p17:sec:prediction}

\noindent
Suppose the grammars are matched. Suppose two subjects have done the labour of translation, read each other's shadows faithfully, and extended to each other the hospitality the first cause requires. They can still fall out of phase---and now by a mechanism that lies not in the code but in the very act by which each subject holds the other in mind. This second cause is, in the author's view, the deepest of the internal ones, and it is the layer at which the present paper makes its first substantive addition to the series. It concerns prediction: what each subject must do to remain coupled to a being it cannot enter, and what happens when that prediction begins, irreparably, to fail.

\subsection{The predictive coupling}
\label{p17:sub:prediction-coupling}

A subject cannot read another's interior directly; opacity, as the introduction insisted, is here ontological and not merely technical (\xref{p17:sub:intro-twolevels}). What a subject has instead is a \emph{model}---a generative model of the other, built up over the history of the bond, by which it anticipates what the other will do, say, want, and mean, and against which it registers what the other actually does. The series has used this predictive-coding vocabulary before, in its account of how trust is \emph{built}: trust, on that account, is what accrues when a subject's prior model of the other is borne out, when self-exposure lowers the cost of being wrong, and when an active leap of inference is rewarded by the other's response.\footnote{See \textcite{paperxiii}, the chapter unifying the three channels of trust under predictive coding: the structural or prior channel, the self-exposure or likelihood channel, and the active-inference channel. The present section runs the same mechanism toward the opposite outcome.} The coupling of two subjects, in this language, is the ongoing mutual tuning of two generative models, each predicting and being predicted, each updating on the error between expectation and event.

The good case is a coupling in which prediction error, when it arises, is \emph{informative}---attributable to a definite source, and so usable to update the model toward a better fit. The other did something unexpected; the unexpected thing has a legible cause---a mood, a circumstance, a thing not said---and the model, taking the cause on board, predicts a little better next time. Error of this kind is not the enemy of coupling but its very nutrition. A bond grows precisely by being surprised and learning from the surprise.

\subsection{When error cannot be attributed}
\label{p17:sub:prediction-unattrib}

The disturbance enters when prediction error arrives \emph{without an attributable source}. The other deviates from the model, repeatedly, and the deviation cannot be assigned: it is impossible to tell whether the error means \emph{the other has changed}---so that the model is now simply out of date and must be revised; or whether it is \emph{noise in the channel}---a passing perturbation signifying nothing, to be ignored on pain of overfitting; or whether \emph{one's own prior was wrong all along}---so that the model was never a good model of this other, and what is collapsing is not the bond but an illusion about it. These three attributions call for opposite responses---update, ignore, or demolish---and the catastrophe of this layer is that the subject \emph{cannot tell which is owed}. The error is real and persistent; its meaning is undecidable.

\begin{claim}[The unattributability of error]
What triggers the bond's self-protection is not prediction error as such---error is the nutrition of coupling---but error that \emph{cannot be attributed}. When a subject cannot decide whether persistent deviation means the other has changed, the channel is merely noisy, or its own prior was mistaken, no update is safe: to revise may be to chase noise; to ignore may be to miss a real change; to demolish may be to destroy a sound bond over a transient. It is the \emph{undecidability of attribution}, not the magnitude of error, that does the damage.
\end{claim}

\subsection{Active withdrawal: estrangement as defensive disinvestment}
\label{p17:sub:prediction-withdrawal}

Faced with persistent, unattributable error and no safe update, a predictive system has one further move, and it is the move that constitutes this layer's form of estrangement. It can reduce its \emph{coupling}---lower the gain on its predictions of the other, narrow the range over which it stakes anything on being right, cease to expose itself to the error by ceasing, in some measure, to predict the other at all. In the vocabulary of active inference, the system can act so as to bring its sensory stream back within the range its model can already handle, not by improving the model but by \emph{withdrawing from the part of the world that the model cannot track}. Applied to a bond, this means: to stop reaching, to lower the stakes, to invest less of oneself in anticipating the other---to pass, henceforth, at a safe distance.

This is the second cause's central and, I think, original claim: that \emph{surechigai} at this layer is not coldness, not indifference, and not the failure of love, but the \emph{defensive disinvestment of a predictive system that can no longer attribute its errors}. The withdrawal is, on its own terms, rational---a way of stanching a loss that cannot be stopped by any available update. And it is, precisely because it is rational and self-protective, almost impossible to read correctly from the other side, where it presents not as self-protection but as the cooling it is so often mistaken for. One subject withdraws to stop bleeding; the other reads the withdrawal as the withdrawal of love, and---being also a predictive system facing now its own unattributable error---withdraws in turn. The structure of this spiral is the matter of the section on power proper (\xref{p17:sec:ir}); here it is enough to have named its engine.

\subsection{The same mechanism, the opposite outcome}
\label{p17:sub:prediction-symmetry}

It must be made explicit what has happened, because it is the section's contribution to the architecture of the series. The mechanism that produces estrangement here is \emph{the very same mechanism} that the earlier paper showed to produce trust. Trust was the active-inference leap rewarded: a subject exposed itself, predicted the other generously, and found the prediction borne out, and the coupling tightened. Estrangement is that identical machinery running toward the opposite pole: a subject exposed itself, predicted the other, found the prediction defeated by errors it could not attribute, and the coupling, in self-defence, loosened. There are not two mechanisms, one benign and one malign. There is one predictive coupling, and its direction---toward tightening or toward withdrawal---is set by whether the errors it encounters can be attributed and so metabolised, or cannot.

\begin{claim}[The reversibility of the predictive mechanism]
Trust-formation and estrangement are not distinct processes but a single predictive coupling traversed in opposite directions. The same machinery that tightens a bond when its errors are attributable loosens it when they are not. This identity has a practical consequence reserved for the praxis (\xref{p17:sec:praxis}): the work that protects a bond at this layer is not the elimination of error---which would be the elimination of coupling itself---but the maintenance of the conditions under which error remains \emph{attributable}.
\end{claim}

\noindent
That consequence belongs to the later argument, but its shape is already visible, and it confirms the paper's governing thesis from an unexpected angle. One cannot, at this layer, act directly on the love---cannot will the coupling to tighten, cannot legislate trust into being. One can act only on the \emph{conditions of attribution}: on whether the channel is clear enough, the exposure mutual enough, the shared history legible enough, that errors arriving between the two can be assigned a source and so turned back into the nutrition of the bond rather than the trigger of its retreat. The political economy of conditions, announced in the introduction as the limit of what the symbolic can reach, here receives its first concrete instance. And it points already beyond the dyad: for whether error can be attributed depends not only on the two subjects but on the field around them---on whether third parties, scripts, and powers are injecting noise into the channel or clarity into it. With that, the internal causes have been laid out as far as the dyad alone can carry them, and the inquiry must widen to the field. But one internal layer remains: the layer at which the failure is not of code, and not of prediction, but of the value that the coupling was supposed to generate---the layer at which the phase itself stops accumulating.
